% resumo em português
\begin{resumo}
\noindent

Este trabalho investigou técnicas computacionais aplicadas ao ensino de matemática por meio de uma revisão sistemática da literatura. O relato foi estruturado com apoio do PRISMA 2020 e identificou 9.431 registros em quatro fontes científicas, dos quais 17 estudos foram incluídos após deduplicação, triagem e avaliação de elegibilidade. Os estudos incluídos foram examinados pelo MMAT 2018, mantendo-se separadas as respostas a cada critério metodológico, sem cálculo de pontuação geral ou ranking de qualidade. A síntese indicou predominância de abordagens supervisionadas voltadas à predição de desempenho e à estimativa de proficiência, além de lacunas relacionadas à explicabilidade, à integração curricular, à participação docente, à reprodutibilidade e à validação em contextos escolares diversos. A fundamentação também mostrou que desempenho observado, proficiência estimada, competência e aprendizagem não são conceitos equivalentes; por isso, resultados computacionais devem apoiar, e não substituir, a interpretação do professor. Como próxima etapa, o trabalho avaliará tecnologias, fontes de dados e alternativas de prototipação com base nos requisitos e limitações identificados.

\vspace{0.2cm}
Palavras-chave: ensino de matemática; aprendizado de máquina; revisão sistemática; avaliação de competências; tecnologia educacional.

\end{resumo}

% resumo em inglês
\begin{resumo}[Abstract]
\begin{otherlanguage*}{english}
\noindent

This work investigated computational techniques applied to mathematics education through a systematic literature review. Reporting was structured with support from PRISMA 2020 and identified 9,431 records across four scientific sources, of which 17 studies were included after deduplication, screening, and eligibility assessment. The included studies were appraised with MMAT 2018 by preserving the response to each methodological criterion, without calculating an overall score or quality ranking. The synthesis showed a predominance of supervised approaches aimed at performance prediction and proficiency estimation, as well as gaps involving explainability, curriculum integration, teacher participation, reproducibility, and validation across diverse school contexts. The theoretical foundation also showed that observed performance, estimated proficiency, competence, and learning are not equivalent concepts; therefore, computational outputs should support, rather than replace, teacher interpretation. As the next stage, the work will evaluate technologies, data sources, and prototyping alternatives based on the identified requirements and limitations.

\vspace{0.2cm}
Keywords: mathematics education; machine learning; systematic review; competency assessment; educational technology.
\end{otherlanguage*}
\end{resumo}
